\documentclass[sigconf]{acmart}

\usepackage{tabularx}

\AtBeginDocument{%
  \providecommand\BibTeX{{%
    Bib\TeX}}}

\setcopyright{none}
\copyrightyear{2026}
\acmYear{2026}
\acmDOI{}
\acmConference[ISWC Workshop 2026]{ISWC Workshop 2026}{2026}{}
\acmISBN{}
\settopmatter{printacmref=false}
\renewcommand\footnotetextcopyrightpermission[1]{}

\begin{document}

\title[CampusClaw]{CampusClaw: Human-in-the-Loop Agentic Actions for Explainable Campus Matching}

\author{Zhicheng Xu}
\affiliation{%
  \institution{Nankai University}
  \city{Tianjin}
  \country{China}}
\email{xzc@mail.nankai.edu.cn}

\author{Zhen Yang}
\affiliation{%
  \institution{Nankai University}
  \city{Tianjin}
  \country{China}}
\email{yz@mail.nankai.edu.cn}

\author{Kaixuan Wen}
\affiliation{%
  \institution{Nankai University}
  \city{Tianjin}
  \country{China}}
\email{wkx@mail.nankai.edu.cn}

\author{Zeshuai Zhu}
\affiliation{%
  \institution{Nankai University}
  \city{Tianjin}
  \country{China}}
\email{zzs@mail.nankai.edu.cn}

\author{Houxin Chang}
\affiliation{%
  \institution{Nankai University}
  \city{Tianjin}
  \country{China}}
\email{chx@mail.nankai.edu.cn}

\renewcommand{\shortauthors}{Xu et al.}

\begin{abstract}
Campus communities depend on fragmented channels such as group chats, bulletin boards, informal marketplaces, and social feeds to support team formation, peer-to-peer trading, dating, and information exchange. Large language models (LLMs) can reduce the friction of form filling, search, and first-contact communication, but directly connecting an LLM to application APIs may introduce unsafe execution, hallucinated recommendations, and unclear user responsibility. This paper presents CampusClaw, a working campus application that explores a human-in-the-loop agentic interaction pattern for localized matching. CampusClaw converts natural-language requests into schema-grounded action proposals, validates them through a backend skill registry, and lets users choose whether to execute, only fill a draft, or only keep the generated suggestion. Its recommendation module returns only real database objects and attaches scores and explanations for team recruitment, dating, and second-hand trading. We describe the design process, system architecture, action proposal pipeline, explainable matching logic, and preliminary user feedback. CampusClaw demonstrates how LLM-based assistance can be embedded into campus applications without giving the model unconstrained authority over user data or application state.
\end{abstract}

\begin{CCSXML}
<ccs2012>
 <concept>
  <concept_id>10003120.10003121.10003122.10003334</concept_id>
  <concept_desc>Human-centered computing~User studies</concept_desc>
  <concept_significance>500</concept_significance>
 </concept>
 <concept>
  <concept_id>10003120.10003121.10003124.10010866</concept_id>
  <concept_desc>Human-centered computing~User interface design</concept_desc>
  <concept_significance>500</concept_significance>
 </concept>
 <concept>
  <concept_id>10010147.10010257.10010293.10010294</concept_id>
  <concept_desc>Computing methodologies~Natural language generation</concept_desc>
  <concept_significance>300</concept_significance>
 </concept>
 <concept>
  <concept_id>10002951.10003317.10003338</concept_id>
  <concept_desc>Information systems~Retrieval models and ranking</concept_desc>
  <concept_significance>300</concept_significance>
 </concept>
</ccs2012>
\end{CCSXML}

\ccsdesc[500]{Human-centered computing~User studies}
\ccsdesc[500]{Human-centered computing~User interface design}
\ccsdesc[300]{Computing methodologies~Natural language generation}
\ccsdesc[300]{Information systems~Retrieval models and ranking}

\keywords{large language models, human-in-the-loop interaction, explainable recommendation, campus computing, agentic applications}

\maketitle

\section{Introduction}
University students frequently need to find teammates, trade second-hand items, ask for campus information, or initiate social connections. In practice, these needs are scattered across WeChat groups, forums, campus marketplaces, and informal social channels. The result is not simply an absence of information, but a mismatch between user intent and the fragmented places where useful opportunities appear. Students must repeatedly search, filter, compare, and decide how to start a conversation.

LLMs provide a promising interface for such scenarios. A student can express a goal in natural language, e.g., ``recommend several frontend teammates for an AI competition'' or ``help me write a polite first message to contact the seller.'' However, deploying LLMs inside interactive social applications raises three risks. First, an LLM should not be allowed to freely call arbitrary backend APIs. Second, a generative model may fabricate users, posts, or products if recommendations are produced as free text. Third, campus matching is socially sensitive: users often want assistance, but they still want final control over what is published, sent, or saved.

We address these issues through \textbf{CampusClaw}, a campus-oriented social and collaboration platform. The implemented application contains forum posting, team recruitment, dating matching, trade matching, messaging, profile management, reporting, blocking, and an AI assistant. The key design idea is to treat the LLM as a \emph{proposal generator} rather than an autonomous operator. The model helps interpret user intent and draft structured payloads, while the backend validates each action against a whitelist of skills and the user decides whether to execute the action.

The research question behind the system is: \emph{how can a campus application use LLMs to reduce search and communication friction while preserving structured grounding, user agency, and operational safety?} We use the term \emph{schema-grounded} to emphasize that the model output is not accepted as an executable instruction by itself. It must be mapped to a typed application intent, constrained by allowed fields, connected to existing database entities where relevant, and confirmed by the user. This framing turns the LLM from an opaque executor into a bounded component within a larger socio-technical workflow.

This paper contributes:
\begin{itemize}
    \item A schema-grounded action proposal pipeline that turns natural language into validated, user-confirmed application actions.
    \item An explainable recommendation mechanism that only recommends real database objects and presents scores and reasons.
    \item A human-in-the-loop interaction design with three choices: fill and publish, fill only, or generate only.
    \item A working campus application prototype that integrates LLM assistance into multiple concrete scenarios rather than a standalone chatbot.
\end{itemize}

\section{Design Motivation and Requirements}
The design of CampusClaw follows a human-centered process rather than starting from model capability alone. In the project presentation and prototype study, we focused on four recurring campus connection scenarios: finding team members, finding dating or interest-based matches, trading second-hand goods, and browsing campus forum information.

\subsection{User Research}
We conducted a small competition-oriented user research process with students aged 18--25 who had experience with project teams, campus activities, second-hand trading, or campus forums. The study combined survey questions, informal interviews, and task-based prototype use. The product presentation reports four pain points: fragmented information channels, time-consuming filtering, discomfort when initiating contact, and missed time-sensitive opportunities. In the preliminary survey summary, 82\% of respondents reported information fragmentation, 76\% reported filtering fatigue, 68\% reported embarrassment or pressure when contacting strangers, and 71\% reported missing opportunities because information arrived too late.

These values are not intended as a large-scale statistical claim. Instead, they motivated the product direction: a system should not only aggregate campus content, but also explain why a connection is relevant and reduce the pressure of first contact.

\subsection{Requirements}
The observations led to four requirements.

\textbf{R1: Unified opportunity discovery.} The application should aggregate forum posts, team recruitment, dating candidates, and trade posts into a coherent product flow.

\textbf{R2: Explainable matching.} A recommendation should not only output a target, but also explain why it is relevant, e.g., shared skills, matched interests, or keyword overlap.

\textbf{R3: Human-controlled automation.} AI should reduce writing and search effort, but final publication or message sending should remain under user control.

\textbf{R4: Safe backend execution.} The LLM should not directly construct URLs or freely invoke backend operations. All executable actions should be routed through a backend registry with validation.

These requirements are relevant for semantic web and knowledge-enabled systems: the core challenge is not merely generating fluent text, but grounding agentic behavior in structured entities, typed actions, and explainable relationships.

\begin{table}[t]
\centering
\caption{Design principles derived from the user research and system risks.}
\label{tab:principles}
\begin{tabularx}{\linewidth}{lX}
\toprule
\textbf{Principle} & \textbf{Operationalization in CampusClaw}\\
\midrule
Groundedness & Recommendations must refer to existing users, posts, teams, or trade items.\\
User agency & Publishing, messaging, and interest signals require explicit confirmation.\\
Progressive automation & Users can choose execution, draft filling, or text-only generation.\\
Context preservation & The assistant receives page context and returns page-specific action proposals.\\
Risk minimization & Destructive actions are excluded from the assistant skill registry.\\
\bottomrule
\end{tabularx}
\end{table}

\section{System Overview}
CampusClaw is implemented as a uni-app/Vue3 client and a Django REST backend. The client can be packaged as a WeChat mini-program or an Android APK. The backend provides modules for accounts, profiles, forum posts, team recruitment, dating profiles and preferences, trade posts, chat messages, moderation, and assistant actions.

\begin{figure*}[t]
\centering
\fbox{\parbox{0.92\textwidth}{
\centering
\vspace{0.9cm}
\textbf{Figure Placeholder: CampusClaw Architecture}\\
Client AI sheet / assistant page $\rightarrow$ assistant API $\rightarrow$ orchestrator $\rightarrow$ BlueLM call and recommendation logic $\rightarrow$ action proposal $\rightarrow$ skill execution after user confirmation.
\vspace{0.9cm}
}}
\caption{CampusClaw uses the LLM as a bounded reasoning component. The backend orchestrator creates proposals, and the skill registry executes only validated actions after explicit user confirmation.}
\Description{A placeholder architecture diagram showing the data flow from the client assistant interface to the backend orchestrator, LLM service, action proposal, and skill execution after user confirmation.}
\label{fig:architecture}
\end{figure*}

The assistant subsystem is organized around four backend components: \texttt{assistant/agent.py} for BlueLM-compatible chat completion calls, \texttt{assistant/orchestrator.py} for turn planning and prompt construction, \texttt{assistant/recommendations.py} for database-grounded recommendation, and \texttt{assistant/skills.py} for executable action registration. On the frontend, \texttt{services/assistant.ts}, \texttt{AssistantSheet.vue}, and \texttt{AssistantActionCard.vue} handle assistant requests, half-screen interaction, and proposal-card rendering. This structure separates model invocation, orchestration, business execution, and client rendering.

\subsection{Data and Control Boundaries}
CampusClaw separates three types of information. First, \emph{conversational context} captures the user's current request and the page from which it was sent. Second, \emph{application entities} represent durable objects such as posts, profiles, applications, trade listings, and chat messages. Third, \emph{action proposals} are temporary records that connect a model-generated suggestion with a possible backend operation. This separation is important because an LLM response can be re-generated or discarded, whereas an application entity changes the persistent social state of the system.

The backend therefore treats assistant output as provisional. A proposal can expire, remain pending, be filled into a form, or be executed after confirmation. The model never receives a raw ability to mutate the database; it only produces structured content that is later interpreted by deterministic code. This boundary is modest but consequential: it shifts safety from prompt wording alone to a combination of prompt constraints, schema validation, entity lookup, permission checks, and user confirmation.

\section{Schema-Grounded Action Proposals}
The core interaction primitive is an \textbf{AssistantActionProposal}. Instead of returning only a natural-language answer, the system may return a structured action card with a type, title, target page, preview, fill payload, status, and expiration time.

\subsection{Pipeline}
The proposal pipeline consists of five steps.

\textbf{Step 1: Intent interpretation.} When a user sends a message, the backend routes it to \texttt{plan\_assistant\_turn()}. The orchestrator first checks whether the prompt is a recommendation or icebreaker request. Otherwise, it asks the LLM to infer the user's intent and missing fields.

\textbf{Step 2: Payload generation.} The orchestrator constructs prompts that include the current page type, context path, recent history, allowed intent names, and allowed fields. The model is asked to return structured JSON rather than an unconstrained command.

\textbf{Step 3: Backend validation.} The returned fields are sanitized and normalized. Only fields allowed for the selected intent are kept. For example, a team recruitment action may include title, summary, details, target size, tags, and required skills; a trade action may include post type, title, description, price, condition, tags, and negotiability.

\textbf{Step 4: User-facing proposal.} The frontend renders an action card. The user can choose ``fill and publish'', ``fill only'', or ``generate only''. This preserves user control and supports low-risk previewing.

\textbf{Step 5: Skill execution.} If the user confirms execution, the backend checks ownership, status, expiration, and skill registration before creating a post, submitting an application, favoriting a trade post, or sending a chat message.

\begin{figure}[t]
\centering
\fbox{\parbox{0.9\linewidth}{
\centering
\vspace{0.8cm}
\textbf{Figure Placeholder: Proposal Lifecycle}\\
Prompt $\rightarrow$ JSON draft $\rightarrow$ validation $\rightarrow$ pending proposal $\rightarrow$ fill / execute / expire
\vspace{0.8cm}
}}
\caption{Lifecycle of an assistant action proposal. The model output is not directly executed; it becomes a temporary object that can be inspected, filled into a form, executed after confirmation, or discarded.}
\Description{A placeholder lifecycle diagram showing that a user prompt becomes a JSON draft, then a validated pending proposal, and finally one of three outcomes: fill, execute, or expire.}
\label{fig:lifecycle}
\end{figure}

\begin{table*}[t]
\centering
\caption{Examples of schema-grounded AI actions in CampusClaw.}
\label{tab:actions}
\begin{tabularx}{\textwidth}{lX X}
\toprule
\textbf{Action Kind} & \textbf{Structured Payload} & \textbf{User-Controlled Outcome}\\
\midrule
Forum post creation & title, body, category, tags & Fill draft or publish after confirmation\\
Forum interaction & post/comment identifier, body or like signal & Comment, reply, or like after confirmation\\
Team recruitment & title, summary, details, target size, skills & Create an open team post or fill an application\\
Trade post creation & type, title, description, price, condition, tags & Create a trade listing or fill a draft\\
Dating signal & target user, interested/skip signal & Express interest or skip after confirmation\\
Context chat message & target user, source type, message body & Send a confirmed private message\\
\bottomrule
\end{tabularx}
\end{table*}

\subsection{Safety Boundary}
The system explicitly rejects a direct tool-use pattern in which the model freely selects arbitrary APIs. The LLM proposes; the backend validates; the user confirms; the registry executes. High-risk actions such as deleting content, blocking users, hiding conversations, or bypassing moderation are not exposed as assistant skills. This does not make the system immune to every prompt-injection risk, but it narrows the executable surface to typed, auditable, and user-confirmed actions.

The design also distinguishes \emph{low-risk content assistance} from \emph{state-changing operations}. Generating an application message, rewriting a post, or suggesting a search query can remain text-only. Creating a post, sending a private message, or expressing interest in another user changes the social state of the system and therefore requires an explicit confirmation step. This distinction mirrors a broader principle for agentic applications: autonomy should increase only when the consequence of error is low, reversible, or clearly reviewable by the user.

\section{Explainable Matching and Recommendation}
CampusClaw includes three recommendation scenarios: team recruitment, dating matching, and trade matching. The system deliberately separates retrieval from generation. Recommended objects must already exist in the database; the LLM is not allowed to invent a user, team, or product.

\textbf{Team recommendation.} The function \texttt{build\_team\_recommendations()} uses the current user prompt as the main signal. It matches prompt terms against team post titles, summaries, details, tags, and required skills. It filters out the user's own posts, closed posts, and posts whose current size has reached the target size. The result includes a recommendation score and a reason such as keyword or skill overlap.

\textbf{Dating recommendation.} The function \texttt{build\_dating\_recommendations()} combines saved dating preferences with the current prompt. Hard preference filters, such as age, height, weight, and gender ranges, are applied unless the user explicitly relaxes them. The current prompt is used as a temporary ranking signal, for example preferring candidates who mention photography or similar interests.

\textbf{Trade recommendation.} The function \texttt{build\_trade\_recommendations()} focuses on item and transaction keywords in the user's prompt. It matches title, description, condition, tags, and trade type, filters out the user's own posts, and ranks open trade posts. The output supports viewing details, favoriting, or contacting the seller.

\begin{table}[t]
\centering
\caption{Recommendation grounding policies.}
\label{tab:recommendations}
\begin{tabularx}{\linewidth}{lX}
\toprule
\textbf{Scenario} & \textbf{Grounding and Explanation}\\
\midrule
Team & Real open posts; explanation based on skills, tags, and text overlap\\
Dating & Visible profiles; explanation based on saved preferences and current prompt\\
Trade & Open trade posts; explanation based on item keywords and trade metadata\\
\bottomrule
\end{tabularx}
\end{table}

The recommendation card includes a title, subtitle, score, score label, reason, and next-step buttons such as contact, fill application, favorite, interested, or skip. This makes the system closer to explainable semantic matching than to a black-box chatbot response. It also creates a bridge to future semantic web extensions: users, posts, skills, interests, trade items, and actions can be modeled as typed entities and relations rather than unstructured strings.

\subsection{Entity-Centric Grounding}
The grounding mechanism can be viewed as a lightweight entity graph even though the current implementation uses relational tables. Users author posts, posts carry tags, team posts request skills and contain vacancies, dating profiles expose bounded attributes and interests, trade posts describe items and transaction states, and messages connect two users. The assistant operates over this typed substrate by retrieving candidate rows, computing overlaps, and generating explanations from candidate attributes.

\begin{table}[t]
\centering
\caption{Core entities used for recommendation and action grounding.}
\label{tab:entities}
\begin{tabularx}{\linewidth}{lX}
\toprule
\textbf{Entity} & \textbf{Grounding role}\\
\midrule
User/Profile & Identity, Claw ID, public profile fields, ownership checks\\
Forum Post & Content discovery, comments, replies, likes, image context\\
Team Post & Skills, target size, current members, application state\\
Dating Profile & Visible attributes, interests, preference compatibility\\
Trade Post & Item title, price, condition, status, favorite/contact actions\\
Action Proposal & Temporary bridge between model output and execution\\
\bottomrule
\end{tabularx}
\end{table}

This structure reduces a common failure mode of generative recommendation: the assistant cannot simply invent a compelling candidate. A recommendation card must be backed by a row that can be opened, contacted, favorited, or acted upon. The generated explanation is therefore not the source of truth; it is a human-readable interpretation of structured evidence.

\section{Human-in-the-Loop Interaction}
CampusClaw's assistant is accessible through a half-screen sheet and an independent assistant page. The half-screen design allows the user to ask for help while staying on the current page. This is important for contextual tasks: a user viewing a team post may ask for an application message, while a user viewing a trade post may ask for an icebreaker message.

\begin{figure}[t]
\centering
\fbox{\parbox{0.9\linewidth}{
\centering
\vspace{0.8cm}
\textbf{Figure Placeholder: AI Action Card}\\
Buttons: Fill and publish / Fill only / Generate only
\vspace{0.8cm}
}}
\caption{The action card exposes automation levels explicitly. Users may execute the backend action, fill a draft for later editing, or keep the generated text only.}
\Description{A placeholder figure showing an AI action card with three user choices: fill and publish, fill only, and generate only.}
\label{fig:hitl}
\end{figure}

The three user choices represent different levels of automation:
\begin{itemize}
    \item \textbf{Fill and publish:} the user confirms the preview and the backend executes the corresponding skill.
    \item \textbf{Fill only:} the generated payload is stored as a pending draft and loaded into the target form, allowing further editing.
    \item \textbf{Generate only:} the assistant reply remains as text, with no change to application state.
\end{itemize}

This design is especially useful for socially sensitive interactions. Users may accept an AI-generated icebreaker only after seeing the full text. The model can reduce anxiety, but the user still decides whether to send. In the dating scenario, low-pressure actions such as ``interested'' and ``skip'' are separated from direct chat initiation. In team and trade scenarios, contact suggestions and draft messages are shown before execution.

\subsection{Interaction Design Rationale}
The half-screen assistant was chosen to preserve spatial continuity with the user's current task. A full-screen chatbot tends to detach the user from the object being discussed, whereas a contextual sheet allows the user to ask for help while still seeing the post, match, or trade item. This is particularly important for explanations: a recommendation reason is more useful when the relevant card remains visible and actionable.

The three-button action card also addresses different levels of trust. Novice users can choose ``generate only'' and manually copy or ignore the suggestion. Users who trust the draft but want editorial control can choose ``fill only''. Users who understand the preview and want efficiency can choose ``fill and publish''. In this sense, the interface does not assume a single ideal degree of automation; it lets the user calibrate automation at the moment of action.

\section{Preliminary Evaluation}
CampusClaw is currently positioned as a working competition prototype. We therefore report preliminary, small-scale feedback rather than a formal deployment study.

The prototype evaluation involved task-based use across three representative flows: (1) ask AI to recommend team posts and submit a team application, (2) ask AI to recommend a trade item and contact the seller, and (3) use AI to generate a first-contact message. We observed whether users could complete the task, whether they understood why an item was recommended, and whether they preferred direct execution or draft filling.

The product presentation reports a small usability test with 28 participants. In that internal test, AI-generated icebreaker messages were associated with a higher reply rate in prototype scenarios, and AI-filled application drafts were associated with higher application completion. Because the test was small and competition-oriented, we treat these numbers as early evidence for design usefulness, not as a generalizable causal result.

Qualitative feedback was consistent with the design requirements. Users reported that recommendation reasons made matches feel less random, team application cards were clearer than searching through chat groups, and AI-generated opening messages reduced the difficulty of starting a conversation. During iteration, we fixed concrete interaction issues including input layout, search result display, contact buttons, image rendering on real devices, and clearer Chinese wording.

\begin{table}[t]
\centering
\caption{Evaluation protocol used for the prototype study.}
\label{tab:evaluation}
\begin{tabularx}{\linewidth}{lX}
\toprule
\textbf{Task} & \textbf{Observed evidence}\\
\midrule
Team recommendation & Whether the user understood recommendation reasons and could submit or fill an application.\\
Trade recommendation & Whether the user could identify a relevant item and choose contact or favorite actions.\\
Icebreaker generation & Whether the generated opening message was considered usable, editable, or inappropriate.\\
Action confirmation & Whether users noticed that execution required explicit confirmation.\\
\bottomrule
\end{tabularx}
\end{table}

The evaluation also revealed an important design tension. Users appreciated automation when it reduced repetitive writing or surfaced relevant opportunities, but they were less comfortable when the system appeared to act on their behalf. This supports the decision to make action proposals visible and revocable rather than silently executing inferred intentions. For future studies, we plan to separate perceived usefulness, perceived control, trust in explanations, and task completion time as distinct outcome measures.

\section{Discussion}
CampusClaw illustrates a middle path between traditional GUI applications and fully autonomous agents. It does not replace structured application logic with free-form model behavior. Instead, it places LLM output inside a typed, inspectable, and revocable proposal workflow.

For semantic web and knowledge-enabled systems, this architecture suggests two directions. First, campus entities such as users, skills, posts, interests, trade items, and actions can be represented as a lightweight domain schema. Second, future versions may replace simple keyword and attribute overlap with ontology-aware or knowledge-graph reasoning, e.g., inferring that ``data visualization'' relates to D3.js, dashboard design, and frontend development.

There are limitations. The current prototype uses relational database rows and hand-written matching functions rather than a formal ontology. The evaluation is small-scale and competition-oriented. Production deployment would require stronger security policies, privacy review, HTTPS deployment, domain registration, and larger-scale field testing. The current system should therefore be read as a design and engineering probe into safe agentic campus applications, not as a finished campus-wide infrastructure.

\subsection{Risks and Mitigations}
Agentic interfaces introduce risks that are not fully solved by better prompts. A prompt may be ambiguous, a user may ask for an unsafe shortcut, or a model may produce a plausible but incomplete payload. CampusClaw addresses these issues at several layers: limiting the skill registry, filtering generated fields, checking object ownership, expiring proposals, and requiring visible confirmation for state-changing actions.

\begin{table}[t]
\centering
\caption{Risk analysis for agentic campus interactions.}
\label{tab:risks}
\begin{tabularx}{\linewidth}{lX}
\toprule
\textbf{Risk} & \textbf{Mitigation}\\
\midrule
Hallucinated target & Recommend only database-backed objects.\\
Unauthorized execution & Validate proposal ownership and pending status.\\
Prompt injection & Expose only whitelisted non-destructive skills.\\
Social overreach & Require confirmation before messaging or interest signals.\\
Loss of user control & Provide generate-only and fill-only alternatives.\\
\bottomrule
\end{tabularx}
\end{table}

These mitigations are intentionally pragmatic. They do not claim to provide a complete formal safety proof. Instead, they create a layered defense suitable for a prototype in which the consequences of model error should remain inspectable and recoverable.

\subsection{Implications for Agent-Native Applications}
CampusClaw suggests that near-term agent-native applications may not remove graphical interfaces. Rather, they may coexist with conventional interfaces while gradually shifting frequent, repetitive, or ambiguous tasks into structured assistant flows. In this transitional mode, buttons remain available for transparency and fallback, while the assistant organizes drafts, recommendations, explanations, and next actions.

The most important design challenge is therefore not whether an LLM can produce text, but how an application can translate uncertain language into reliable state transitions. Schema-grounded proposals are one answer: they make the boundary between language and execution explicit. This boundary can later be enriched with ontologies, provenance metadata, policy constraints, and more rigorous explanation models.

\section{Ethics and Privacy Statement}
CampusClaw is designed for socially sensitive campus scenarios, including dating, private messaging, and peer-to-peer trading. The prototype therefore avoids exposing high-risk assistant actions such as deleting user content, blocking other users, or bypassing moderation. Before production deployment, the system would require institutional privacy review, clearer consent flows, stronger identity and abuse-prevention mechanisms, and a public policy for data retention and moderation.

\section{Conclusion}
We presented CampusClaw, a campus matching and collaboration application that integrates LLM assistance through schema-grounded, human-in-the-loop action proposals. The system uses BlueLM to understand intent, generate drafts, recommend real database objects, and produce contact suggestions, while backend validation and user confirmation prevent unconstrained execution. The prototype demonstrates how agentic interaction can be embedded into a practical campus application without turning the LLM into an unsafe autonomous operator. Future work will focus on stronger semantic modeling, knowledge-graph-based matching, and larger-scale user evaluation.

\begin{thebibliography}{10}
\bibitem{brown2008design}
Brown, T.: Design thinking. Harvard Business Review 86(6), 84--92 (2008).

\bibitem{berners2001semantic}
Berners-Lee, T., Hendler, J., Lassila, O.: The Semantic Web. Scientific American 284(5), 34--43 (2001).

\bibitem{hogan2021knowledge}
Hogan, A., Blomqvist, E., Cochez, M., d'Amato, C., de Melo, G., Gutierrez, C., Kirrane, S., Gayo, J.E.L., Navigli, R., Neumaier, S., et al.: Knowledge Graphs. ACM Computing Surveys 54(4), 1--37 (2021).

\bibitem{floridi2018ai}
Floridi, L., Cowls, J., Beltrametti, M., Chatila, R., Chazerand, P., Dignum, V., Luetge, C., Madelin, R., Pagallo, U., Rossi, F., et al.: AI4People---An ethical framework for a good AI society. Minds and Machines 28(4), 689--707 (2018).

\bibitem{xi2023agents}
Xi, Z., Chen, W., Guo, X., He, W., Ding, Y., Hong, B., Zhang, M., Wang, J., Jin, S., Zhou, E., et al.: The rise and potential of large language model based agents: A survey. arXiv preprint arXiv:2309.07864 (2023).

\bibitem{qin2023tool}
Qin, Y., Liang, S., Ye, Y., Zhu, K., Lu, C., Lin, Y., Jindal, P., Wang, X., Yuan, Z., Sanwal, G., et al.: Tool learning with foundation models. arXiv preprint arXiv:2304.08354 (2023).

\bibitem{zanzotto2020hitl}
Zanzotto, F.M.: Human-in-the-loop artificial intelligence. Wiley Interdisciplinary Reviews: Data Mining and Knowledge Discovery 10(6), e1382 (2020).
\end{thebibliography}

\end{document}
